\section{Limitations}

Several limitations constrain the interpretation of this study. First, MLL23 is the single internal dataset and AML-LMU is the single external domain. Additional independent datasets would be needed to estimate the range of domain-shift behavior across scanners, stains, laboratories, annotation policies, and patient populations. Second, reliable patient or acquisition-group identifiers are unavailable in the working MLL23 manifests. The V2 split is a conservative exact and near-duplicate sensitivity analysis, not patient-level validation. AML-LMU documents subject counts in the source collection, but patient identifiers were not available in the audited fields used for this project.

Third, the unknown classes are defined by the protocol. They are outside the known training taxonomy, but they do not define malignancy, leukemia, abnormal diagnosis, or a positive clinical finding. Fourth, the class distribution is imbalanced in both internal and external data. This motivates macro-F1, balanced accuracy, per-morphology results, and open-set metrics, but imbalance still affects uncertainty around small external unknown classes such as atypical lymphocytes, metamyelocytes, monoblasts, and smudge cells.

Fifth, the model family is limited to ImageNet-pretrained ResNet18 variants with frozen training settings. The study uses five seeds, which is stronger than a single-run report but not exhaustive. Sixth, ArcFace is evaluated with one frozen margin and scale. Further tuning could change results, but tuning after seeing confirmatory outcomes would reopen method selection and weaken the frozen evaluation design. Seventh, the internal unknown set was inspected during prior analysis, so internal explanatory findings should not be treated as independent discovery. The external test-only analysis is therefore especially important.

Eighth, no external adaptation is performed. This strengthens the external validation claim but does not estimate the best possible performance after domain adaptation, stain normalization, external calibration, or target-domain fine-tuning. Ninth, persistent-homology conclusions are limited to the evaluated filtrations, vectorizers, fusion designs, sample sizes, and summary statistics. They should not be generalized to all topological methods. Finally, the study is a research evaluation on public datasets, not a prospective clinical workflow evaluation, not a patient-level diagnostic task, and not a guarantee that dataset-defined unknowns cover the full diversity of real-world hematological morphology.
